\section{涨落耗散定理}

\subsection{响应与平衡涨落}

由 Kubo 公式 \eqref{eq:kubo}，响应函数的虚部编码耗散。另一方面，平衡态算符涨落的频谱由动态结构因子（或对称关联函数）给出。涨落耗散定理（FDT）把二者联系起来。

对一般算符 $A$、$B$，在频率空间（平衡态）
\begin{equation}
  \mathrm{Im}\,\chi^R_{BA}(\omega)
  =\frac{1}{2\hbar}
  \bigl(1-e^{-\beta\hbar\omega}\bigr)
  S_{BA}(\omega),
\label{eq:FDT_gen}
\end{equation}
其中 $S_{BA}$ 为相应涨落谱（具体预因子随 $S$ 的 $2\pi$ 约定略有不同）。密度通道的形式见 \eqref{eq:FDT_chi}。

\subsection{推导要点}

\begin{enumerate}
  \item 用 Lehmann 表示同时写出 $\chi^R$ 与 $S$；
  \item 矩阵元相同，能量 $\delta$ 函数相同，差别仅在于热因子与对易子相对编序期望的组合；
  \item 利用 $e^{-\beta E_n}/e^{-\beta E_m}=e^{-\beta(E_n-E_m)}$ 整理即得 FDT。
\end{enumerate}

\subsection{物理含义}

\begin{itemize}
  \item 测量耗散（吸收）等价于探测相应算符的平衡涨落；
  \item 经典高温极限 $\hbar\omega\ll k_BT$ 时，$\mathrm{Im}\,\chi\sim(\omega/2k_BT)S_{\mathrm{sym}}$；
  \item 导电、磁化率、密度响应均可纳入同一框架。
\end{itemize}

\begin{note}
FDT 是平衡统计与线性响应之间的桥梁。电子液中，由 $S(q,\omega)$ 经 \eqref{eq:FDT_chi} 得 $\mathrm{Im}\,\chi$，再经 Kramer--Kronig 得实部，从而构造 $\varepsilon(q,\omega)$。
\end{note}
